零多项式的情形立即成立。对非零数据，归一化为 \(\|f\|_2=1\)，并针对给定环带固定孔径 \(\gamma\) 及常数 \(K\)，见 \ref{ms:prop-5-1}。先设 \(h<1/(8K)\)。令 \(a=Kh\)、\(b=2a\)，用 \(E(T,\tau)=\mathcal F(T,\tau)\) 表示命题~\ref{ms:prop-5-1} 的能量。该命题给出能量在开环域 \(a<|\tau|<4\) 上的非负性与光滑性，该环域包含闭环域 \(b\le|\tau|\le3\)。它还给出 \(\Delta_{T,\tau}E\le V(|\tau|)E\) 在 \(b<|\tau|<3\) 上成立，其中
\(V(s)=Cs^{-2}e^{-cs/h}\)。
对于 \ref{ms:lemma-6-1}，所需积分一致有界，因为
\[
 \int_b^3s\log(s/b)V(s)\,ds
 \le C\int_1^\infty \frac{\log t}{t}\,e^{-2cKt}\,dt .
\]
另有 \(V\) 在 \([1/2,3/2]\) 上一致有界。\ref{ms:prop-5-1} 在两个圆周上的界各自在整条圆周上一致。因而 圆周平均 \(H\)（对 \(E\) 取平均） 满足
\[
 H(0,2b)\le Ca^2\|f\|_{p_c}^{p_c},\qquad
 H(T,b)\ge ca^2\|U(T)f\|_{p_c}^{p_c}\quad(1\le T\le2).
\]
用 \ref{ms:lemma-6-1}，取 \(r=2b\)，约去 \(a^2\)，并取 \(p_c\) 次方根。由于 \(K\) 固定且 \(\log(3/(2Kh))\lesssim\log(2/h)\)，得所述断言。
所得不等式是齐次的，可通过重新伸缩去掉归一化。当 \(h\ge1/(8K)\) 时，所有有贡献的频率都在固定的球 \(|\xi|\le8KC_0\) 内。由 Hölder 不等式，每个这种频率处的 Fourier 系数不超过 \(\|f\|_{p_c}\) 的常数倍，且 \(e^{-iT|\xi|}\) 的模为一。对所得有限个逐点上界求和，即得此范围内的估计。
